\section{Evaluation under imbalance and dependence}
\label{sec:bg:eval}

A \emph{sample} is one market at one moment at which a prediction is asked
(\cref{sec:bg:notation}); a validation day has about ten million of them.  The three
classes are of very unequal size: at one hour more than half of the samples are
\textsc{const}, at one minute nine in ten.  Plain accuracy, the share of samples predicted
correctly, would therefore reward a model that always predicts \textsc{const}: it is right
89\% of the time at one minute and knows nothing.  The thesis uses \emph{balanced
accuracy} instead \cite{brodersen2010balanced}.  For each class, the \emph{recall} is the
share of the samples of that class that the model predicts as that class, how many of
the real \textsc{up} samples it called \textsc{up}, and so on; balanced accuracy is the
mean of the three recalls,
\begin{equation}
  \mathrm{BA} = \frac{1}{3} \sum_{c=1}^{3} \mathrm{recall}_c, \qquad
  \mathrm{recall}_c = \frac{\#\{i : y_i = c,\ \hat{y}_i = c\}}{\#\{i : y_i = c\}} .
\end{equation}
A model that always predicts one class has recall one on that class and zero on the other
two, hence balanced accuracy $1/3$ whatever the class sizes, and random guessing scores
the same.  Balanced accuracy depends on the predicted classes, not on the probabilities,
which is why the decision rule of \cref{eq:logit-adjustment} matters as much as the model
and why every comparison of the thesis is made after the same calibration.

A confidence interval says how much a number would move if the experiment were repeated
on fresh data.  The usual interval treats every sample as an independent observation, and
with ten million samples it would be a few hundredths of a point wide.  Here it would be
wrong, because the samples are far from independent: one market contributes thousands of
samples a few seconds apart that are almost copies of one another, and all the markets of
one hour share the same news, so there are far fewer independent observations than
samples.  The intervals are therefore computed with a \emph{block bootstrap}
\cite{kunsch1989jackknife}.  The validation day is cut into its 24 hours; a new ``day''
is built by drawing 24 hours at random, with replacement, from those 24; the metric, or
the difference between two models, is recomputed on it; this is repeated two thousand
times, and the 2.5th and 97.5th percentiles of the two thousand values are the 95\%
interval.  Drawing whole hours keeps the dependence inside an hour intact, which is what
makes the interval account for it.  A difference whose interval includes zero is reported
as a tie.  Every number is also given on the \emph{deduplicated} subset of
\cref{sec:data:labels}, one sample per market per fifteen minutes, on which the samples
are closest to independent, and on the \emph{extreme} and \emph{central} price ranges
separately; a result that holds on one view only is reported as such.
